\section{Construcción y verificación del modelo: detalle}
\label{anx:modelo}
% Origen: 05-paso4-construccion-verificacion.tex y los anexos de parámetros y de la fila CSV, movidos al anexo.

\subsection{Herramienta y estructura}
\label{anx:estructura}

El modelo \texttt{SubteConstitucion.alp} se construyó en AnyLogic~8.9 en su edición gratuita (PLE). Tiene
dos capas independientes dentro del mismo archivo:

\begin{itemize}
  \item \textbf{Capa lógica (\texttt{Main}).} Usa la Process Modeling Library y el circuito
        \texttt{Source} $\to$ \texttt{Queue} $\to$ \texttt{Seize} $\to$ \texttt{Delay} $\to$
        \texttt{Release} $\to$ \texttt{Sink}, con los molinetes como un \texttt{ResourcePool}. Cada oleada
        se inyecta completa con \texttt{inject(n)}. Implementa los escenarios \EO{} a E3 con una cola común
        y servidores intercambiables. Fue el primer prototipo, y se conserva como referencia verificada y
        como base de E2 y E3.
  \item \textbf{Capa peatonal (\texttt{MainPeatonal}).} Usa la Pedestrian Library y el circuito
        \texttt{PedSource} $\to$ \texttt{PedService} $\to$ \texttt{PedGoTo} $\to$ \texttt{PedSink}. Los
        peatones caminan por el vestíbulo, eligen molinete, esperan en su cola, validan y salen por una
        línea física al otro lado de la batería. Es la capa que produce los resultados del estudio.
\end{itemize}

Las capas están separadas para que un cambio gráfico en el vestíbulo no altere las métricas de la capa
lógica.

\subsection{Recorrido del pasajero}
\label{anx:recorrido}

El modelo sigue al pasajero del Roca desde que baja del tren hasta que valida en un molinete y entra a la zona
paga. Hay dos niveles peatonales, con el molde del ejemplo oficial \emph{Stairs} de AnyLogic. De fondo, la vista
muestra el plano de la estación enviado por Trenes Argentinos y las capas de OpenStreetMap, georreferenciados de
modo que el pasillo del modelo coincide con el real (anexo~\ref{anx:geometria}).

\begin{itemize}
  \item \textbf{Andén (nivel 0).} Cada tren se detiene en una de las siete vías con andén, la primera libre. Su
        pasaje baja por las 21 puertas de la formación (7 coches de 23~m con 3 puertas por lado) y camina el
        andén hasta la cabecera. La dispersión con que la oleada llega al hall sale de esa caminata.
  \item \textbf{Hall (nivel 0).} El pasajero cruza el hall entre los obstáculos dibujados con la silueta del
        plano: las cabeceras de los andenes, el bloque central de escaleras, los vestuarios y las boleterías.
  \item \textbf{Escaleras mecánicas.} Baja por uno de los dos grupos de escaleras mecánicas del bloque central:
        el norte si bajó en un andén al norte del centro del bloque y el sur en otro caso. Cada escalera es un
        área en pendiente de $10{,}5 \times 2{,}8$~m, con 4,6~m de desnivel y velocidad máxima de 0,685~m/s, de
        modo que la bajada lleva $4{,}6/0{,}30 \approx 15{,}3$~s. La bajada se camina y puede congestionarse. Las
        escaleras fijas del centro del bloque no se modelan.
  \item \textbf{Subsuelo (nivel $-4{,}60$~m).} Al pie de las escaleras está la plaza de transbordo, de donde sale
        un pasillo en L de 84~m y 5~m de ancho sobre la traza del pasillo subterráneo de OpenStreetMap. La plaza
        y el ancho del pasillo son supuestos.
  \item \textbf{Vestíbulo.} Mide $50 \times 34$~m. La calle entra por la escalera B3 de la calle Brasil, en su
        posición real. La zona no paga tiene boletería, máquinas de carga y columnas como obstáculos. Los
        molinetes forman dos bancos de 14 pasos de 1~m separados por una cabina, y la zona paga empieza donde el
        pasillo real se conecta con el andén de la Línea~C.
  \item \textbf{Molinetes.} Son servicios lineales de 1,6~m que se atraviesan de la zona no paga a la paga. Cada
        uno tiene su propia cola de 6~m con la cabeza a 2~m del paso, como en el ejemplo oficial \emph{Subway
        Entrance Hall}.
\end{itemize}

El interior del vestíbulo es supuesto: los 34~m de ancho son los que hacen falta para alojar 28 molinetes. Los 22
primeros representan los del vestíbulo Principal y los 6 restantes, la ampliación hipotética E1-b. Los escenarios
no duplican la geometría: \EO{} suspende los molinetes 21 a 28 y \EI{} los 23 a 28, con la API de
\texttt{ServiceWithLine}. La escalera B4 no se modela, porque baja del otro lado de la conexión con el andén y,
sin el plano de SBASE, no se sabe qué molinetes usan sus pasajeros.

\begin{figure}[H]
\centering
\figuraSiExiste{tpi-subte-modelo-espacial.png}{\textwidth}
\caption{Modelo espacial en AnyLogic: andenes del Roca, hall con el bloque central de escaleras, plaza de
transbordo, pasillo y vestíbulo, sobre el plano de Trenes Argentinos.}
\label{fig:plano}
\end{figure}

\subsection{Presentación del modelo}
\label{anx:presentacion}

La vista del modelo se organizó en cuatro bloques, según la consigna de la cátedra, para que alguien que no lo
conoce pueda leerlo de arriba abajo (figura~\ref{fig:bloques}):
\begin{enumerate}
  \item \textbf{Variables}: parámetros del escenario y de la corrida, datos embebidos, generadores aleatorios
        y supuestos, cada uno con una línea que dice qué es;
  \item \textbf{Medidas de rendimiento}: las primarias, las secundarias, las de validación y un panel en vivo
        con el día, el escenario, la hora, el P90, la cola y la gente en tránsito;
  \item \textbf{Lógica}: el diagrama de flujo del pasajero, los bloques de los trenes, los eventos y las
        funciones;
  \item \textbf{Modelo espacial}: la animación en 2D y 3D, con dos mapas de calor (hall y subsuelo) que
        acumulan la densidad máxima de la corrida y marcan en rojo las zonas por encima de
        1,5~personas/m².
\end{enumerate}

\begin{figure}[H]
\centering
\figuraSiExiste{tpi-subte-modelo-bloques.png}{\textwidth}
\caption{Vista del modelo en cuatro bloques: variables, medidas de rendimiento y lógica arriba; modelo espacial
abajo.}
\label{fig:bloques}
\end{figure}

\subsection{Lógica de la capa peatonal}
\label{anx:logica}

\textbf{Demanda.} Cada réplica es un día hábil real, sorteado con la semilla entre los 75 días de marzo a junio de
2026 con servicio normal del Roca (sección~\ref{sec:paso2-datos-modelo-conceptual}). Del día se toman dos cosas,
ambas embebidas en el modelo para que funcione igual en AnyLogic Cloud:
\begin{itemize}
  \item las validaciones del vestíbulo Principal por ventana de 15 minutos, de SBASE;
  \item la llegada real de cada tren a Plaza Constitución, al minuto, de Trenes Argentinos.
\end{itemize}
Cada tren lleva la demanda acumulada desde el tren anterior de su ramal. El pasaje concreto es una variable de
Poisson con esa media, con su propio generador. El total del Roca se reescala para que sea el 80~\% de las
validaciones del día. La demanda de la calle se genera por adelgazamiento (\emph{thinning}) de un proceso de
Poisson con la tasa de cada ventana. Si falta un parámetro sin dato, la inicialización se detiene con un error
explícito en lugar de usar un valor oculto.

\textbf{Elección de cola.} \texttt{PedService} delega la elección en una función propia,
\texttt{elegirColaPed}. Para cada molinete habilitado, la función calcula el desvío lateral a 1,3~m/s más
los pasajeros ya asignados a ese molinete que todavía no empezaron a validar, multiplicados por el tiempo
de validación de referencia. El contador incluye a los que caminan hacia la cola y no solo a los que ya
están en ella. Así se evita el \enquote{efecto manada}: una oleada completa eligiendo la misma cola vacía.

\textbf{Medición.} La espera de cada pasajero empieza en el evento \texttt{onEnterQueue} y termina al
comenzar la validación. Las métricas de la franja se congelan a las 09:30 ($t = 9000$), pero la corrida
continúa hasta que sale el último pasajero ingresado antes de esa hora. Recién entonces el modelo escribe
una fila de 20 campos en el archivo de corridas (anexo~\ref{anx:csv}) y, en otro archivo, las validaciones
simuladas y observadas por ventana del mismo día. Antes de escribir exige que se cumpla
inyectados $=$ generados $=$ procesados $+$ atascados. Si no se cumple, se detiene con un error de
conservación. Los atascados son los pocos pasajeros que el cierre por atasco excluye al final de la corrida
(anexo~\ref{anx:pilotos-produccion}); en una corrida normal valen cero.

\textbf{Ruteo.} Un elemento Pedestrian Settings fija el peso de densidad del ruteo en cero: cada peatón va por
el camino más corto en lugar de esquivar las zonas cargadas. No cambia ninguna decisión del modelo, porque la
escalera la determina el lado del hall y el molinete la regla de elección, y vale igual para los cuatro
escenarios. Se adoptó por costo de cómputo (anexo~\ref{anx:pilotos-produccion}) y se declara como simplificación.

\textbf{Números aleatorios comunes.} Cada fuente de azar tiene su propio generador, sembrado a partir de la
semilla de la réplica: el día, la carga de los trenes, la calle, el tiempo de validación y la velocidad de
marcha. Los cuatro escenarios de una réplica reciben el mismo día, los mismos trenes, los mismos pasajeros en
los mismos instantes, con las mismas velocidades \parencite[cap.~11]{law2015}. Lo único que cambia es la
cantidad de molinetes habilitados y la distribución del tiempo de validación.

\textbf{Automatización.} El experimento de variación de parámetros \texttt{PeatonalCorridasApareadas}
ejecuta 40 corridas con evaluaciones en paralelo. La corrida $j$ usa la semilla $20260923 + \lfloor j/4 \rfloor$
y el escenario $j \bmod 4$: \EO{}, \EI{}, \ET{} y \EIT{}. Cada grupo de cuatro corridas consecutivas comparte la
semilla y, por lo tanto, el día. Un script (\texttt{cargar\_corridas.py}) valida las filas y las transfiere a la
planilla de análisis sin transcripción manual. Rechaza filas de demostración, corridas incompletas, réplicas sin
\EO{}, escenarios con distinto día o distinta demanda y violaciones de conservación.

\subsection{Verificación}
\label{anx:verificacion}

La verificación comprueba que el programa hace lo que el modelo conceptual especifica. Se hizo en tres
niveles.

\textbf{Verificador automático.} El script \texttt{verificar\_modelo.py} revisa el archivo del modelo sin
modificarlo:
\begin{itemize}
  \item la estructura del XML, la unicidad de los identificadores y las referencias entre elementos;
  \item los parámetros de cada experimento;
  \item que el horario cargado en el modelo sea exactamente el extraído de los PDF oficiales;
  \item que los 75 días embebidos coincidan con los datos de SBASE y de Trenes Argentinos, y que en ninguno de
        ellos dos trenes ocupen la misma vía a la vez;
  \item la orientación de las 28 colas, la geometría de las escaleras y de la plaza, y que cada elemento de la
        vista quede dentro de su bloque.
\end{itemize}
Además compila las funciones Java del modelo contra las bibliotecas instaladas de AnyLogic. Las ejercita
en un arnés de prueba con casos de resultado conocido:
\begin{itemize}
  \item \EO{} suspende 8 molinetes y \EI{}, 6;
  \item una oleada se reparte entre colas en lugar de amontonarse en una;
  \item quien encuentra el molinete libre tiene espera cero y no se le suma la caminata;
  \item no se generan ingresos a partir de las 09:30;
  \item la proporción de pasajeros que llega en tren y la que llega desde la calle coinciden con las
        configuradas;
  \item los cuatro escenarios reciben los mismos ingresos y sortean los mismos tiempos de validación, pasajero
        por pasajero;
  \item la media empírica del servicio triangular coincide con la teórica;
  \item la fila de salida tiene los 20 campos en el orden de la planilla.
\end{itemize}
También comprueba que la planilla calcula los mismos intervalos de confianza que una implementación
independiente en Python, con una diferencia menor que $10^{-6}$. La última ejecución, del 10/11/2026,
pasó todos los controles.

\textbf{Traza manual.} En la capa lógica se reprodujo una oleada de 5 pasajeros con 2 molinetes y 3~s de
servicio, cuyo resultado se calcula a mano. Las esperas son 0, 0, 3, 3 y 6~s: espera media de 2,4~s, P90
de 6~s, cola máxima de 3, área bajo la cola de 12 pasajero-segundo y ocupación de 15 molinete-segundo. El
modelo da exactamente esos valores.

\textbf{Controles en cada corrida.} Cada réplica de producción comprueba la conservación y que haya
drenado por completo. Si la corrida termina sin vaciar el sistema, emite una fila marcada como
\texttt{INCOMPLETO} que el cargador rechaza.

\subsection{Control de versiones}
\label{anx:versiones}

El modelo, los scripts, los datos procesados y este informe se guardan en un repositorio Git. Para evitar
que el IDE y un editor externo se pisen, el archivo del modelo tiene un solo editor por vez. Al cerrar el
IDE se revisa la diferencia con la versión anterior, se corre el verificador y se registra el cambio. Las
salidas de cada corrida se versionan junto al modelo que las produjo.

\subsection{Parámetros del modelo en la producción}
\label{anx:parametros}

\begin{table}[H]
\centering
\caption{Parámetros de \texttt{MainPeatonal} en \texttt{PeatonalCorridasApareadas}}
\label{tab:anx-parametros}
\small
\begin{tabularx}{\textwidth}{@{}>{\raggedright\arraybackslash}p{4.4cm} l X@{}}
\toprule
Parámetro & Valor & Significado \\
\midrule
\texttt{molinetesOperativosPed} & 20 o 22 & Molinetes habilitados \\
\texttt{servicioPedSeg} (min, máx) & 2,4 (1,8; 3,5) o 1,8 (1,5; 2,5) & Validación triangular: SUBE o sin contacto \\
\texttt{proporcionRocaPed} & 0,80 & Parte de las validaciones que llega en tren [supuesto] \\
\texttt{demoraAccesoRocaSeg} & 10~s & Del arribo a la apertura de puertas [supuesto] \\
\texttt{duracionDescargaSeg} & 30~s & Descenso por las 21 puertas [supuesto] \\
\texttt{diaSbasePed} & $-2$ & Día sorteado con la semilla entre los 75 días \\
\texttt{llegadasRealesPed} & verdadero & Trenes con la llegada real del día \\
\texttt{cargaTrenPoissonPed} & verdadero & Pasaje de cada tren con distribución de Poisson \\
\texttt{atrasoMaxRocaSeg} & 0 & Atraso adicional (solo para sensibilidad con el horario oficial) \\
Velocidad de marcha & U(1,1; 1,5)~m/s & Generador propio \\
Escaleras mecánicas & 0,685~m/s & 10,5~m con 4,6~m de desnivel a 0,30~m/s vertical \\
\bottomrule
\end{tabularx}
\end{table}

\subsection{Esquema de la fila CSV\_PEATONAL}
\label{anx:csv}
% Origen: filaResultadoPed() del modelo y CAMPOS de cargar_corridas.py (20 campos).

Cada corrida escribe una fila separada por punto y coma. El primer campo es el prefijo:
\texttt{CSV\_PEATONAL} (cargable), \texttt{CSV\_PEATONAL\_DEMO} (demostración, nunca se carga) o
\texttt{CSV\_PEATONAL\_INCOMPLETO} (corrida sin drenar, bloquea la carga).

\begin{longtable}{@{}r>{\raggedright\arraybackslash}p{5.2cm}>{\raggedright\arraybackslash}p{7.5cm}@{}}
\toprule
Campo & Nombre & Detalle \\
\midrule
\endhead
1 & Escenario & \EO{}, \EI{}, \ET{} o \EIT{} \\
2 & Semilla & Común a los cuatro escenarios de la réplica \\
3 & Pasajeros generados & Igual en los cuatro escenarios de la réplica \\
4 & Pasajeros desviados & Cero en todos los escenarios evaluados \\
5 & Procesados a las 09:30 & Validaciones hasta el corte \\
6 & Procesados con drenaje & Igual a los generados (conservación) \\
7 & Espera media & Segundos, desde la entrada a la cola hasta el inicio de la validación \\
8 & P90 de espera & Primaria \\
9 & Proporción con espera mayor a 30~s & Primaria, fracción entre 0 y 1 \\
10 & Lq & Largo medio de la cola en la franja \\
11 & Cola máxima & Personas \\
12 & Tiempo de drenaje final & Último egreso menos último ingreso \\
13 & Utilización media & Fracción, entre los molinetes habilitados \\
14 & Dispersión de utilización & Desvío entre molinetes \\
15 & Pasajeros de la cohorte pico & Ingresados de 08:15 a 08:45 \\
16 & Espera media de la cohorte pico & Segundos \\
17 & P90 de la cohorte pico & Primaria \\
18 & Día SBASE & Índice del día sorteado (0 a 74) \\
19 & P90 de tránsito del Roca a la cola & Segundos desde que baja del tren hasta que entra a la cola \\
20 & CV de validaciones entre molinetes & Para validar el reparto contra SBASE \\
\bottomrule
\end{longtable}
